\documentclass[10pt,a4paper]{amsart}
\usepackage[margin=19mm]{geometry}
\usepackage{amsmath,amssymb,mathtools}
\usepackage[scr=rsfs,cal=euler]{mathalfa}
\usepackage{xfrac}
\usepackage[unicode,hidelinks,bookmarks=false]{hyperref}
\newcommand{\ie}{i.e.\ }
\newcommand{\cf}{cf.\ }
\newcommand{\resp}{respectively\ }
\newcommand{\Subsection}{\subsection}
\providecommand{\nobreakdash}{\nobreak\hbox{-}}
\setlength{\emergencystretch}{2em}
\multlinegap0pt
\usepackage{siunitx}
\usepackage{siunitx}
\usepackage{siunitx}
\usepackage{siunitx}
\usepackage{xcolor}
\usepackage{siunitx}
\usepackage{amssymb}
\usepackage{mathtools}
\newtheorem{theorem}{Theorem}
\newcommand{\NN}{\mathbb N}
\newcommand{\ind}[1]{\mathbf{1}_{\{#1\}}}
\title{Optimal sequential tests yield log-optimal e-processes}
\author{Ashwin Ram, Aaditya Ramdas}
\date{Source: arXiv:2605.12720v1 (CC BY 4.0)}
\begin{document}
\maketitle

\noindent\emph{Source and licence.} Optimal sequential tests yield log-optimal e-processes. Version arXiv:2605.12720v1 (CC BY 4.0), distributed by arXiv under the Creative Commons Attribution 4.0 International licence, http://creativecommons.org/licenses/by/4.0/, as recorded in the arXiv licence metadata for this exact version and shown on the abstract page https://arxiv.org/abs/2605.12720v1. Full licence notice: \texttt{licenses/arxiv-2605.12720-CC-BY-4.0.txt}.\par\smallskip

\noindent\textit{Editorial scope.} Counted unit: the article's theorem thm:main (source0.tex lines 268--279). The article's complete original proof of it is delivered with it (source0.tex lines 629--698, one \texttt{proof} environment of 70 lines, which the article places away from the statement). No mathematics is written, repaired or re-derived: the statement and the proof are the article's own lines, in the article's own notation, and no retained line is substituted or re-flowed.  Retained uncounted as the source-local context the proof needs: lines 196--198; lines 205--207; lines 226--228; lines 231--233; lines 236--238.  Nothing was omitted from any retained span, so the package carries no omission record.  No retained line contains a citation, so the wrapper carries no bibliography and no bibliography entry.  No retained line contains a figure environment, an image-inclusion command or drawing code, so no image is delivered and the package declares no asset dependency.  Editorial formatting only, in addition: an AMS \texttt{amsart} preamble that restates the article's own declarations and macros used by the retained lines, namely the environments \texttt{theorem} on separate counters, together with the standard packages those lines need.  The retained environments print in this document as Theorem 1 in order of appearance, whereas the article prints the same items with the numbering of its own full document.  The complete native source of the article as distributed in the arXiv e-print for this exact version is bundled unchanged as \texttt{sources/200440/source0.tex}.
\begin{equation}\label{eq:V}
\sup_{P\in\mathcal P} P(\tau_k<\infty)\le \alpha_k.
\end{equation}

\begin{equation}\label{eq:AO}
\frac{\tau_k}{b_k}\longrightarrow \frac{1}{I},
\end{equation}

\begin{equation}\label{eq:W}
\sum_{k=1}^{\infty}w_k\alpha_k\le 1.    
\end{equation}

\begin{equation}\label{eq:aggregate}
M_t:=\sum_{k=1}^\infty w_k \ind{\tau_k\le t}.
\end{equation}

\begin{equation}\label{eq:Wprofile}
W(x):=\sum_{k:b_k\le x}w_k.    
\end{equation}

\begin{theorem}\label{thm:main}
Assume that \eqref{eq:V}, \eqref{eq:W}, and \eqref{eq:AO} hold. Suppose that $M$ is the aggregate (as in \eqref{eq:aggregate}) and let $W$ be the profile (as in \eqref{eq:Wprofile}). Suppose in addition $W(x)\to\infty$ and for some $\rho\in[0,1]$ we have that
    \begin{equation}\label{eq:R}
    \frac{\log W(x)}{x}\longrightarrow \rho.   
    \end{equation}
    Then, almost surely under $Q$ as $t\to\infty$,
    \begin{equation}\label{eq:G}
    \boxed{
    \frac{1}{t}\log M_t\longrightarrow \rho I.  
    }
    \end{equation}
\end{theorem}

\begin{proof}[Proof of Theorem~\ref{thm:main}]
In proving this theorem, we will first prove the following two statements.

\begin{enumerate}
    \item For every $\varepsilon\in(0,1)$ there exists an almost sure event $\Omega_\varepsilon$ such that for every $\omega\in\Omega_\varepsilon$ there is an integer $K_\varepsilon(\omega)$ for which for all $k\ge K_\varepsilon(\omega)$,
    \begin{equation}\label{eq:sandwich-tests}
    \frac{1-\varepsilon}{I}b_k\le \tau_k(\omega)\le \frac{1+\varepsilon}{I}b_k. 
    \end{equation}
    \item Take some $\varepsilon\in(0,1)$ and $\omega\in\Omega_\varepsilon$. Set $C_\varepsilon(\omega):=\sum_{k=1}^{K_\varepsilon(\omega)-1}w_k$. Then it follows that for every $t\in\mathbb N_0$,
    \begin{equation}\label{eq:M-sandwich}
    W\left(\frac{It}{1+\varepsilon}\right)-C_\varepsilon(\omega)\le M_t(\omega) \le W\left(\frac{It}{1-\varepsilon}\right)+C_\varepsilon(\omega).    
    \end{equation}
\end{enumerate}

We will first prove \textcolor{red}{statement (1)}. Consider some $\varepsilon\in(0,1)$. Thanks to \eqref{eq:AO}, the random variables $\tau_k/b_k$ converge $Q$-almost surely to $1/I$. Therefore the event
\[
\Omega_\varepsilon:=\left\{\omega : \exists K_\varepsilon(\omega)\in\NN \text{ such that } \left| \frac{\tau_k(\omega)}{b_k} - \frac{1}{I} \right| \le \frac{\varepsilon}{I} \text{ for all } k\ge K_\varepsilon(\omega) \right\}
\]
has $Q$-probability one. If we now rearrange the inequality $\left|\frac{\tau_k(\omega)}{b_k}-\frac{1}{I}\right|\le \frac{\varepsilon}{I}$, we have \eqref{eq:sandwich-tests}, proving the claim.

We will now prove \textcolor{red}{statement (2)}. Take some $\varepsilon\in(0,1)$ and $\omega\in\Omega_{\varepsilon}$. With minor abuse of notation, let us write $K:=K_{\varepsilon}(\omega)$ and $C:=C_{\varepsilon}(\omega)=\sum_{k=1}^{K-1}w_k$. We will prove the lower and upper bounds in \eqref{eq:M-sandwich} separately. We will start with the \underline{lower bound}. Consider any $k\ge K$ such that $b_k\le \frac{It}{1+\varepsilon}$. Thanks to \eqref{eq:sandwich-tests}, it follows that
\[
\tau_k(\omega)\le \frac{1+\varepsilon}{I}b_k\le \frac{1+\varepsilon}{I}\cdot\frac{It}{1+\varepsilon} = t.
\]
Thus, it follows that $\ind{\tau_k(\omega)\le t}=1$. If we then sum $w_k$ over all such $k$ we have that
\[
M_t(\omega) \ge \sum_{\substack{k\ge K\\ b_k \le It/(1+\varepsilon)}} w_k.
\]
We will now split the entire profile $W(It/(1+\varepsilon))$ into indices below $K$ and indices at least $K$. That is,
\[
W\left(\frac{It}{1+\varepsilon}\right) = \sum_{\substack{k< K\\ b_k \le It/(1+\varepsilon)}} w_k + \sum_{\substack{k\ge K\\ b_k \le It/(1+\varepsilon)}} w_k.
\]
Therefore,
\[
\sum_{\substack{k\ge K\\ b_k \le It/(1+\varepsilon)}} w_k = W\left(\frac{It}{1+\varepsilon}\right)- \sum_{\substack{k< K\\ b_k \le It/(1+\varepsilon)}} w_k \ge W\left(\frac{It}{1+\varepsilon}\right) - C.
\]
This holds because the subtracted sum will be at most $\sum_{k<K}w_k=C$. Therefore, $M_t(\omega)\ge W\left(\frac{It}{1+\varepsilon}\right)-C$. We will now prove the \underline{upper bound}. To this end, suppose that $k\ge K$ and $\tau_k(\omega)\le t$. By the lower half of \eqref{eq:sandwich-tests}, we have that
\[
\frac{1-\varepsilon}{I}b_k\le \tau_k(\omega)\le t.
\]
If we multiply by $I/(1-\varepsilon)$, we have that $b_k\le \frac{It}{1-\varepsilon}$. Therefore, each index $k\ge K$ that contributes to $M_t(\omega)$ must satisfy $b_k\le It/(1-\varepsilon)$. Therefore, we have
\[
M_t(\omega) \le \sum_{k< K} w_k + \sum_{\substack{k\ge K\\ b_k \le It/(1-\varepsilon)}} w_k \le C + W\left(\frac{It}{1-\varepsilon}\right).
\]
This proves \eqref{eq:M-sandwich}. We will now prove the \textcolor{red}{main statement} of the theorem. Assume to this end that \eqref{eq:R} holds and that $W(x)\to\infty$. Consider some $\varepsilon\in(0,1)$ and let us work on the same full probability event $\Omega_{\varepsilon}$ from the proof of \textcolor{red}{(1)}. Once again, we will write $C=C_{\varepsilon}(\omega)$. Since $W(x)\to\infty$, there must exist some $t_0(\omega)$ such that for every $t\ge t_0(\omega)$, both $W\left(\frac{It}{1+\varepsilon}\right)>2C$ and $W\left(\frac{It}{1-\varepsilon}\right)>C$. Now, for such a $t$, thanks to the lower bound in \eqref{eq:M-sandwich}, we have
\[
M_t(\omega)\ge W\left(\frac{It}{1+\varepsilon}\right)-C\ge \frac{1}{2}W\left(\frac{It}{1+\varepsilon}\right).
\]
On a similar vein, the upper bound implies that
\[
M_t(\omega)\le W\left(\frac{It}{1-\varepsilon}\right)+C\le 2W\left(\frac{It}{1-\varepsilon}\right).
\]
Let us now take the logarithms and divide by $t$. We then have for all $t$ sufficiently large enough that
\[
\frac{1}{t}\log W\left(\frac{It}{1+\varepsilon}\right) - \frac{\log 2}{t} \le \frac{1}{t}\log M_t(\omega) \le \frac{1}{t}\log W\left(\frac{It}{1-\varepsilon}\right) + \frac{\log 2}{t}.
\]
Let us use \eqref{eq:R}. For the lower bound, we have
\[
\frac{1}{t}\log W\left(\frac{It}{1+\varepsilon}\right) = \frac{I}{1+\varepsilon} \cdot \frac{1}{\,It/(1+\varepsilon)\,} \log W\left(\frac{It}{1+\varepsilon}\right) \longrightarrow \frac{\rho I}{1+\varepsilon}.
\]
Similarly, 
\[
\frac{1}{t}\log W\left(\frac{It}{1-\varepsilon}\right)\longrightarrow \frac{\rho I}{1-\varepsilon}.
\]
Hence it follows that $\liminf_{t\to\infty} \frac{1}{t}\log M_t(\omega) \ge \frac{\rho I}{1+\varepsilon}$ and $\limsup_{t\to\infty} \frac{1}{t}\log M_t(\omega) \le \frac{\rho I}{1-\varepsilon}$. This will hold for all $\varepsilon\in(0,1)$ on the corresponding almost sure event $\Omega_{\varepsilon}$. We can now intersect over the countable set $\varepsilon\in \mathbb  Q\cap(0,1)$ to get a new event of $Q$-probability one on which both inequalities hold simultaneously for each rational $\varepsilon\in(0,1)$. We can now let $\varepsilon\downarrow 0$ through the rationals to give us
\[
\liminf_{t\to\infty}\frac{1}{t}\log M_t(\omega)\ge \rho I,\quad \limsup_{t\to\infty}\frac{1}{t}\log M_t(\omega)\le \rho I.
\]
The limit therefore exists and equals $\rho I$. This proves \eqref{eq:G}, completing the proof.
\end{proof}

\end{document}
